% 第 6 章骨架：导引框已定，正文待写。写作计划见 docs/PLAN.md“第 6 章”，写法见 docs/STYLE.md。
\chapter{美索不达米亚与埃及：神庙、宫廷与记账}\label{ch:mesopotamia}

\begin{chapterguide}
\guidequestion{先有物物交换，还是先有记账与债务？一个没有铸币的经济怎样定价、放贷和征税？}
\begin{guidegoals}
\item 描述神庙与宫廷经济中白银与大麦作为计价单位的运作方式；
\item 比较“物物交换起源说”与“信用起源说”，并说明各自依据的史料；
\item 解释利率上限与周期性债务豁免的经济逻辑；
\item 指出这些史料能说明、和不能说明的宏观问题。
\end{guidegoals}
\guideprereq{\chref{ch:money}\enspace\nameref{ch:money}}
\end{chapterguide}

\chapterstub
